\section{Geomonitoring}

\subsection{TPS-Monitoring}
Ein weiterer Teil des 'Vertiefungsprofil Geosensorik und Monitoring II' ist das automatische Monitoring mit Leica MS60 Tachymetern.
Dazu wird das Instrument über ein Python-Skript von einem Raspberry Pi 4 aus angesteuert.
Im Rahmen des IoT-Workshops wurde ein automatisches Monitoring erstellt und vom 03.11.2026 - 04.11.2026 ausgeführt. 


\subsubsection{Einrichten und Aufsetzen}
Um den ersten Umgang mit dem Raspberry Pi zu erlangen wird er zunächst im Schulzimmer eingerichtet und konfiguriert.
Dazu werden alle 'Mini-Computer' auf den neusten Stand gebracht, und den Tachymeter über das Y-Kabel an ihn angeschlossen.
Nachdem zunächst diverse Einstellungen ausprobiert und Tests durchgeführt werden, wird das automatische Monitoring mittels
dem vorgegebenen Python-Skript vorbereitet.
Dies liest lediglich die Messdaten der MS60 aus und publiziert sie auf der Webseite von InfluxDB.
Zudem wird eine Push-Benachrichtigung nach jeder xxx an die Anwendung xxx gesendet. \todo{Push-Benachrichtigung Einfügen}
Die eigentliche Messung wird am Instrument mit dem Satzmessungsprogramm gestartet. Dies wird über die Zeitsteuerung terminiert. 
Da maximal 200 Sätze am Stück möglich sind, wird ein Intervall von 7,5 Minuten gewählt, welcher eine Zeitspanne von 25 Stunden abdeckt.

\begin{figure}[h!]
    \centering
    \includegraphics[width=\textwidth]{Raspy-Pi.jpg}
    \caption{Einrichten des Monitorings mit Tachymeter und Raspberry Pi auf dem Dach des Campus Muttenz.}
    \label{fig:Raspy-Pi-Dach}
\end{figure}


Leider funktionierte der automatische Start auf dem Dach des Campus nicht wie gewünscht, weshalb zu Debugging-Zwecken ein
Bildschirm und Tastatur für den Raspberry Pi angeschafft werden musste. Die Messung konnte um ca. 14:30 Uhr gestartet werden.
Abbildung \ref{fig:Raspy-Pi-Dach} zeigt, dass trotz Kälte die Motivation und der Zusammenhalt in der Gruppe sehr hoch war.

\begin{figure}[h!]
    \centering
    \includegraphics[width=\textwidth]{InfluxDB-Graph.png}
    \caption{Visualisierung der gemessenen Testpunkte in InfluxDB Cloud.}
    \label{fig:Influx-Graph}
\end{figure}

Die ausgelesenen Daten werden über das auf dem Raspberry Pi laufende Python Skript direkt an die InfluxDB Cloud gesendet.
Aus bisher unerklärlichen Gründen erfolgte die online Synchronisation lediglich bis ca. Mitternacht. \todo{Andere Gruppen auch erwähnen?}
Bei einer erneuten Verbindung mit dem Internet am nächsten Morgen konnten auch noch die restlichen Daten in der InfluxDB Cloud gesichert werden.

Wie die Graph-Visualisierung (siehe Abbildung \ref{fig:Influx-Graph}) aus dem Data Explorer der InfluxDB Cloud zeigen,
weisen die in diesem Zeitraum überwachten Punkte keine groben Ausreisser auf. 

\begin{figure}[h!]
    \centering
    \includegraphics[width=\textwidth]{InfluxDB-Histo.png}
    \caption{Visualisierung der Anzahl Messungen (Histogramm) pro Punkt in InfluxDB Cloud.}
    \label{fig:Influx-Histo}
\end{figure}

In der online-Anwendung gibt es einige Darstellungsmöglichkeiten, jedoch sind die weiteren Einstellungen deutlich begrenzt.
Eine genaue Analyse der Daten ist somit nicht möglich.
Ebenfalls ist die Beschreibung der Darstellungsformen unzureichend. 
In der Abbildung \ref{fig:Influx-Histo} wird beispielsweise ein Histogramm dargestellt.
Vermutlich eingefärbt nach der Anzahl Messungen in diesem Zeitraum. Eine Definition dazu ist leider nicht auffindbar.

% GEOCOM-Befehle erwähnen well coool :D


\subsection{Low Cost GNSS-Monitoring}



\FloatBarrier